% Компилировать XeLaTeX (в Overleaf: Menu -> Compiler -> XeLaTeX)
% Опции класса: single | onehalf (по умолчанию) | newpage | bysection
\documentclass{sfureport}

% ---------- данные для титульного листа (форма М / Н СТУ 7.5–07–2021) ----------
\institute{Институт космических и информационных технологий}
\department{Кафедра вычислительной техники}
\worktype{Отчёт по лабораторной работе}   % или: Отчёт о практической работе, Реферат, Эссе, Контрольная работа
%\discipline{информатике}                  % строка «по ...» (нужна для реферата, контрольной, эссе)
\title{Исследование алгоритмов сортировки}
\teacher{Иванов И.\,И.}
\author{Петров П.\,П.}
\studentgroup{КИ00-00Б, 000000000}
\city{Красноярск}
\date{2026}

\begin{document}

\maketitle

% Содержание: по п. 6.4.1 при объёме до 24 страниц включать не рекомендуется.
% Для большой работы раскомментируйте:
% \stutoc

\structsection{Введение}

Настоящий отчёт оформлен по стандарту университета СТУ 7.5–07–2021. Основной текст набран шрифтом Times New Roman размером 14~пт с полуторным межстрочным интервалом, абзацный отступ составляет 12,5~мм, выравнивание~--- по ширине. Поля страницы: левое 30~мм, правое 10~мм, верхнее и нижнее по 20~мм.

Цель работы~--- изучить алгоритмы сортировки и сравнить их по времени выполнения. Для достижения цели необходимо решить следующие задачи:
\begin{itemize}
  \item реализовать алгоритмы сортировки;
  \item измерить время работы на массивах разного размера;
  \item сравнить полученные результаты.
\end{itemize}

\section{Теоретическая часть}

\subsection{Простые алгоритмы}

Простейшие алгоритмы сортировки имеют квадратичную сложность. К ним относятся:
\begin{enumerate}
  \item сортировка пузырьком;
  \item сортировка вставками, в которой различают:
    \begin{enumerate}
      \item обычную вставку;
      \item бинарную вставку;
    \end{enumerate}
  \item сортировка выбором.
\end{enumerate}

Среднее число сравнений при сортировке пузырьком определяют по формуле
\begin{equation}
  C = \frac{n(n-1)}{2},
  \label{eq:cmp}
\end{equation}
\begin{gde}
  $C$ --- число сравнений;
  \gdeitem $n$ --- размер массива.
\end{gde}

Ссылка на формулу: значение $C$ вычисляют по формуле~(\ref{eq:cmp}). Источники указывают так: подробнее см. в работе~\cite{pushkareva}, а конкретную страницу так~\cite[с.~29]{shevelev}.

\subsection{Быстрые алгоритмы}

Быстрые алгоритмы, например быстрая сортировка, в среднем работают за время $O(n \log n)$.

\section{Практическая часть}

Результаты измерений приведены в таблице~\ref{tab:time}, а общая схема эксперимента показана на рисунке~\ref{fig:scheme}.

\begin{table}
  \caption{Время сортировки массивов}
  \label{tab:time}
  \tableunit{В миллисекундах}
  \begin{tabular}{|C{5cm}|C{3cm}|C{3cm}|}
    \hline
    \textbf{Алгоритм} & \textbf{$n = 10^3$} & \textbf{$n = 10^4$}\\
    \hline\hline
    Пузырёк & 4 & 380\\
    \hline
    Вставки & 2 & 190\\
    \hline
    Быстрая сортировка & 0,1 & 1,4\\
    \hline
  \end{tabular}
\end{table}

\begin{figure}
  \includegraphics[width=0.45\textwidth]{example-image}
  \caption{Схема эксперимента}
  \label{fig:scheme}
\end{figure}

Фрагмент реализации приведён в листинге~\ref{lst:bubble}.

\begin{lstlisting}[language=Python,caption={Сортировка пузырьком},label={lst:bubble}]
def bubble_sort(a):
    # внешний цикл по проходам
    for i in range(len(a)):
        for j in range(len(a) - i - 1):
            if a[j] > a[j + 1]:
                a[j], a[j + 1] = a[j + 1], a[j]
    return a
\end{lstlisting}

\structsection{Заключение}

В ходе работы были реализованы алгоритмы сортировки и выполнено сравнение времени их работы. Быстрая сортировка существенно превосходит квадратичные алгоритмы на больших массивах.

% Список использованных источников (п. 7.10). Записи оформляют по ГОСТ Р 7.0.100,
% примеры~— в приложении У стандарта. Ниже образцы из самого стандарта.
\begin{thebibliography}{99}
  \bibitem{pushkareva} Пушкарева, Т. П. Информатика : учебное пособие / Т. П. Пушкарева ; Сибирский федеральный университет, Политехнический институт. – Красноярск : СФУ, 2021. – 132 с. – ISBN 978-5-7638-4497-9.
  \bibitem{shevelev} Шевелев, А. С. Самоцветы Восточной Сибири : иллюстрированное научное издание / А. С. Шевелев ; научный редактор Д. П. Гладкочуб. – Иркутск : Призма, 2020. – 175 с.
\end{thebibliography}

\stuappendix[обязательное]{Исходный код программы}

Текст приложения. Рисунки, таблицы и формулы в приложении нумеруются с буквой приложения, например «Рисунок А.1».

\begin{figure}
  \includegraphics[width=0.4\textwidth]{example-image-a}
  \caption{Пример рисунка в приложении}
\end{figure}

\end{document}
